%%% BEGIN fig:thr-kinematics
\begin{figure}[htbp]
\centering
\begin{tikzpicture}[fd,
  fc hadron/.style={draw=black!70,line width=1.5pt,line cap=butt},
  fc rem/.style={draw=black!45,line width=0.8pt},
  fc soft/.style={gluon,draw=green!45!black},
  fc lab/.style={font=\small},
  fc note/.style={font=\footnotesize},
  fc hb/.style={hard blob,circle,minimum size=0.56cm,inner sep=0pt,font=\footnotesize},
  fc had/.style={hadron blob,circle,minimum size=0.78cm,inner sep=0pt}]
  % =====================================================================
  % (a) partonic threshold: q qbar -> gamma* + soft radiation
  % =====================================================================
  \begin{scope}
    \coordinate (A) at (-3.1,1.3);  \coordinate (B) at (-3.1,-1.3);  \coordinate (H) at (0,0);
    \draw[fc hadron] (-4.35,1.3) -- (A);
    \draw[fc hadron] (-4.35,-1.3) -- (B);
    \draw[fc rem] ($(A)+(0.3,0.2)$) -- (-1.2,1.95);
    \draw[fc rem] ($(A)+(0.35,0.08)$) -- (-1.2,1.62);
    \draw[fc rem] ($(B)+(0.3,-0.2)$) -- (-1.2,-1.95);
    \draw[fc rem] ($(B)+(0.35,-0.08)$) -- (-1.2,-1.62);
    \draw[quark=0.38] (A) -- (H);
    \draw[quark=0.62] (H) -- (B);
    \draw[fc soft] (-1.18,0.494) -- (-0.3,1.5);
    \draw[fc soft] (-1.18,-0.494) -- (-0.3,-1.5);
    \draw[photon] (H) -- (1.75,0);
    \node[fc had] at (A) {$A$};  \node[fc had] at (B) {$B$};
    \node[fc hb] at (H) {$H$};
    \node[vertex] at (-1.18,0.494) {};  \node[vertex] at (-1.18,-0.494) {};
    \path[mom'={$\xi_AP_A$},mom pos=0.3,mom sep=0.28cm] (A) -- (H);
    \path[mom={$\xi_BP_B$},mom pos=0.3,mom sep=0.28cm] (B) -- (H);
    \path[mom'={$k$},mom pos=0.55,mom sep=0.26cm] (-1.18,0.494) -- (-0.3,1.5);
    \node[fc lab,anchor=south west] at (1.75,0.05) {$\gamma^*(q)$};
    \node[fc lab,above] at (-4.0,1.33) {$P_A$};
    \node[fc lab,below] at (-4.0,-1.33) {$P_B$};
    \node[fc note,anchor=west,text=green!35!black] at (0.1,1.45) {软辐射：$E_X\sim Q(1-z)$};
    \node[fc note,anchor=west] at (0.5,-0.75) {$\hat s=\xi_A\xi_Bs$，$z=Q^2/\hat s\to1$};
    \node[fc note,anchor=west] at (0.5,-1.2) {$q^2=Q^2$};
    \node[fc lab] at (-1.2,-2.65) {(a)\ 部分子阈值：末态只有软辐射};
  \end{scope}
  % =====================================================================
  % (b) refactorization of the hard coefficient: H x S_DY
  % =====================================================================
  \begin{scope}[yshift=-5.15cm]
    \node[fc lab] at (-5.0,0) {$\Delta_{q\bar q}(z)\simeq$};
    % hard function: virtual corrections at the scale Q, no real emission
    \begin{scope}[xshift=-2.55cm]
      \coordinate (L) at (-0.75,0);  \coordinate (R) at (0.75,0);
      \draw[quark] (-1.55,0.85) -- (L);
      \draw[quark] (L) -- (-1.55,-0.85);
      \draw[quark] (R) -- (1.55,0.85);
      \draw[quark] (1.55,-0.85) -- (R);
      \draw[photon] (L) -- (R);
      \node[fc hb] at (L) {$H$};  \node[fc hb] at (R) {$H^\dagger$};
      \draw[cut] (0,-1.05) -- (0,1.05);
      \node[fc note] at (0,-1.45) {$H(Q,\mu)$};
    \end{scope}
    \node[fc lab] at (-0.55,0) {$\times\ Q$};
    % soft function: two past-pointing Wilson lines on each side of the cut
    \begin{scope}[xshift=2.4cm]
      \coordinate (L) at (-1.05,0);  \coordinate (R) at (1.05,0);
      \draw[fc soft] (-1.71,0.57) -- (1.71,-0.57);
      \draw[wilson=0.2] (-2.15,0.95) -- (L);
      \draw[wilson=0.55] (L) -- (-2.15,-0.95);
      \draw[wilson=0.55] (R) -- (2.15,0.95);
      \draw[wilson=0.2] (2.15,-0.95) -- (R);
      \node[opvertex] at (L) {};  \node[opvertex] at (R) {};
      \node[vertex] at (-1.71,0.57) {};  \node[vertex] at (1.71,-0.57) {};
      \draw[cut] (0,-1.05) -- (0,1.05);
      \node[fc note,anchor=east] at (-2.15,0.95) {$n_A$};
      \node[fc note,anchor=east] at (-2.15,-0.95) {$n_B$};
      \node[fc note,anchor=west] at (2.15,0.95) {$n_A$};
      \node[fc note,anchor=west] at (2.15,-0.95) {$n_B$};
      \node[fc note] at (0,-1.45) {$S_{\rm DY}(\omega,\mu)$，$\omega=2E_X\simeq Q(1-z)$};
    \end{scope}
    \node[fc lab,anchor=west] at (5.0,0) {$+\ \cdots$};
    \node[fc lab] at (-1.2,-2.3) {(b)\ 领头幂次下硬系数的重新因子化（“$\cdots$”为$(1-z)$的次领头幂次项）};
  \end{scope}
\end{tikzpicture}
\caption{阈值极限$z=Q^2/\hat s\to1$（\cref{thm:thr-refactorization}）。(a)~来自$f_{q/A}$、$f_{\bar q/B}$的部分子（动量$\xi_AP_A$、$\xi_BP_B$；反夸克线上的箭头为费米子数流向）湮灭为$\gamma^*$，末态除强子残余外只剩能量$E_X\sim Q(1-z)$的软辐射（绿色）。
(b)~在$(1-z)$的领头幂次，硬系数重新因子化为$\Delta_{q\bar q}=H(Q,\mu)\,Q\,S_{\rm DY}(Q(1-z),\mu)$（分布$\delta(1-z)$、$\mathcal D_k$）：$H$只含尺度$Q$的虚修正，软函数$S_{\rm DY}$由每侧两条过去指向的 Wilson 线（沿$n_A$、$n_B$，箭头为色流向）构成，测量部分子质心系中软辐射的能量$\omega=2E_X$。
阈值因子化嵌套在共线因子化$\sigma=\sum f_A\otimes f_B\otimes\Delta$之内，其余项只在$(1-z)$的领头幂次受外层定理控制。}
\label{fig:thr-kinematics}
\end{figure}
%%% END fig:thr-kinematics

%%% BEGIN fig:thr-mellin
\begin{figure}[htbp]
\centering
\begin{tikzpicture}[fd,
  fc ax/.style={lcaxis,gray!60!black},
  fc lab/.style={font=\small},
  fc note/.style={font=\footnotesize},
  fc pole/.pic={\draw[line width=0.9pt,line cap=round] (-0.08,-0.08)--(0.08,0.08) (-0.08,0.08)--(0.08,-0.08);},
  fc contour/.style={draw=blue!65!black,line width=1.15pt,postaction={decorate,decoration={markings,
      mark=at position 0.2 with {\arrow{Stealth[length=6pt,width=5pt]}},
      mark=at position 0.66 with {\arrow{Stealth[length=6pt,width=5pt]}}}}},
  fc cut/.style={draw=black!80,line width=0.9pt,decorate,decoration={zigzag,segment length=4pt,amplitude=1.6pt}}]
  % axes
  \draw[fc ax] (-3.4,0) -- (7.6,0) node[below left,fc lab,text=black] {$\Real N$};
  \draw[fc ax] (0,-2.75) -- (0,2.95) node[left,fc lab,text=black] {$\Imag N$};
  % singularities of the PDF moments (schematic)
  \foreach \x in {-2.6,-1.75,-0.9,0.75} \pic at (\x,0) {fc pole};
  \draw[fc cut] (-3.3,0) -- (-2.68,0);
  \node[fc note,align=center,anchor=north] at (-1.75,-0.22) {部分子分布矩\\$f_{A,N}\,f_{B,N}$的奇点};
  \draw[gray!60!black,line width=0.45pt] (0.75,0.12) -- (1.05,0.62) node[above,fc note,text=black] {最右奇点};
  % Landau branch point and cut of the resummed exponent
  \draw[fc cut] (5.45,0) -- (7.35,0);
  \fill (5.45,0) circle[radius=1.9pt];
  \node[fc note,anchor=north] at (5.45,-0.18) {$N_L$};
  \draw[gray!60!black,line width=0.45pt] (5.6,0.12) -- (5.95,0.68) node[above,fc note,text=black] {Landau 支点};
  % minimal-prescription contour: Re N = c, bending to the left at large |Im N|
  \draw[fc contour] (1.35,-2.6) -- (2.35,-1.45) .. controls (2.75,-1.0) and (2.75,-0.8) .. (2.75,-0.45)
      -- (2.75,0.45) .. controls (2.75,0.8) and (2.75,1.0) .. (2.35,1.45) -- (1.35,2.6);
  \fill (2.75,0) circle[radius=1.5pt];
  \node[fc note,anchor=north west] at (2.78,-0.04) {$c$};
  \node[fc lab,text=blue!55!black,anchor=west] at (2.85,1.25) {$C_{\rm MP}$};
\end{tikzpicture}
\caption{Mellin 反演与 Landau 极点（示意）。反演积分$\int_C\frac{\dd N}{2\pi\ii}\,\tau^{-N}f_{A,N}f_{B,N}\Delta_N$（截面的归一化见正文）的围道$C$自下而上。在固定对数精度下，阈值指数按$\alpha_s\ln\bar N$组织；当有效标度$Q^2/\bar N^2$降到跑动耦合$\alpha_s(q^2)$的 Landau 极点时，$\ln\Delta_N$在正实轴远处的$N_L$出现支点（及向右的割线）；
最小处方\cite{CataniEtAl1996}的围道$C_{\rm MP}$与实轴交于$c$，$c$位于部分子分布矩的奇点与$N_L$之间（其余奇点都在围道左侧，$N_L$及其割线在右侧），在$|\Imag N|$大处向左弯折以改善收敛。处方是定义，不是定理（\cref{rem:thr-prescriptions}）。}
\label{fig:thr-mellin}
\end{figure}
%%% END fig:thr-mellin
